Agradeço primeiramente à minha família, que sempre me ofereceu suporte e apoio incondicional. Sou profundamente grato pela compreensão nos momentos de ausência e pelo incentivo constante que tornou possível a realização deste trabalho. Sem a base sólida que me proporcionaram, este percurso teria sido significativamente mais difícil.

Dirijo também meus agradecimentos aos meus amigos, que, mesmo diante das dificuldades do período acadêmico, sempre estiveram presentes, oferecendo apoio emocional e companheirismo. A presença de cada um foi fundamental para que eu mantivesse o equilíbrio ao longo da minha trajetória.

Agradeço igualmente aos professores que contribuíram para minha formação, transmitindo conhecimento, ética e compromisso. Suas orientações ao longo dos anos foram essenciais para o desenvolvimento das competências que me permitiram elaborar este Trabalho de Conclusão de Curso.

Aos colegas da faculdade, registro minha gratidão pela troca de experiências, pelas discussões construtivas e pela colaboração mútua. Cada interação fortaleceu meu aprendizado e ampliou minha visão profissional.

À instituição de ensino, direciono meu agradecimento a cada profissional que, diariamente e de forma incansável, colabora para o funcionamento e a manutenção deste espaço. Sou grato pelos recursos disponibilizados e pelo ambiente propício ao crescimento intelectual. A qualidade e o compromisso com a educação foram determinantes para a construção deste percurso.

E, por fim, agradeço a mim mesmo: pela persistência, pela resiliência e pela capacidade de enfrentar e superar os desafios que surgiram ao longo da minha jornada. Reconheço o esforço pessoal empregado em cada etapa da minha vida acadêmica e me orgulho da dedicação investida, que me permitiu chegar até aqui — ainda que eu saiba que esta conquista não representa um ponto final, mas apenas mais um passo na minha evolução.